\documentclass[11pt]{article}
\usepackage[margin=1in]{geometry}
\usepackage{amsmath,amssymb,amsthm,mathtools}
\usepackage{graphicx}
\usepackage{microtype}
\usepackage[hidelinks]{hyperref}

\newtheorem{theorem}{Theorem}
\newtheorem{lemma}{Lemma}
\newcommand{\B}{\mathcal B}
\newcommand{\U}{\mathcal U}

\title{Contractive Direct Limits Preserve\\
$L^p$-Operator Algebras}
\author{Automated mathematics run \texttt{fa\_banach\_001}, lane 1}
\date{August 9, 2026}

\begin{document}
\maketitle

\begin{center}
\fbox{\parbox{0.9\textwidth}{\textbf{Status.}
Candidate full solution, likely valid, pending expert review.  For every fixed
$1\le p<\infty$, an arbitrary contractive direct limit of $L^p$-operator
algebras is again an $L^p$-operator algebra.  No complete-contractivity or
injectivity hypothesis on the connecting maps is required.}}
\end{center}

\section{Source question}

Phillips~\cite{PhillipsSource} defines an $L^p$-operator algebra to be a Banach
algebra isometrically isomorphic to a norm-closed subalgebra of
$\B(L^p(X,\mu))$.  Proposition~6.1 of that paper constructs the direct limit
$A=\varinjlim_i A_i$ of any directed system of Banach algebras and contractive
homomorphisms.  If $\varphi_i:A_i\to A$ are the canonical maps, it proves
\begin{equation}\label{eq:limitnorm}
 \|\varphi_i(a)\|
 =\lim_{j\ge i}\|\varphi_{j,i}(a)\|,
 \qquad a\in A_i,
\end{equation}
and that $\bigcup_i\varphi_i(A_i)$ is dense in $A$.

Immediately afterward, on PDF page 35, the paper asks whether $A$ is an
$L^p$-operator algebra when every $A_i$ is one.

\begin{figure}[ht]
  \centering
  \includegraphics[width=\textwidth]{figures/open_problem_crop.png}
  \caption{The source question and its immediate context on page 35 of
  Phillips~\cite{PhillipsSource}.}
\end{figure}

Throughout the result below, $1\le p<\infty$, the range of the source paper.

\section{Proof intuition}

Represent each stage $A_j$ isometrically on some $L^p$-space $E_j$.  A stage
element $a\in A_i$ then has a coherent tail of represented images
$\rho_j(\varphi_{j,i}(a))$ for $j\ge i$.  Choose an ultrafilter which regards
every final segment as large and let this operator tail act on the Banach-space
ultraproduct $(E_j)_\U$.

There are two exact norm facts.  First, the norm of an operator induced by a
bounded family is the ultralimit of the coordinate operator norms.  Second,
the ordinary tail limit of those norms is precisely the direct-limit norm in
\eqref{eq:limitnorm}.  Thus the tail construction is isometric.  It is also
well defined when two stage elements become equal only asymptotically, since
then the tail norm of their difference tends to zero.  The dense algebraic
limit consequently embeds isometrically, and the embedding extends to its
completion.  Finally, ultraproducts of $L^p$-spaces are $L^p$-spaces for finite
$p$.

\section{The theorem}

We first record the elementary operator-ultraproduct norm identity used in the
proof.

\begin{lemma}\label{lem:operatornorm}
Let $(E_j)_{j\in I}$ be nonzero Banach spaces, let $\U$ be an ultrafilter on $I$, and
let $(T_j)_{j\in I}$ be a uniformly bounded family with
$T_j\in\B(E_j)$.  The coordinatewise operator
\[
 T_\U:(E_j)_\U\longrightarrow(E_j)_\U,
 \qquad [(x_j)]\longmapsto[(T_jx_j)],
\]
satisfies
\[
 \|T_\U\|=\lim_{j\to\U}\|T_j\|.
\]
\end{lemma}

\begin{proof}
The upper bound follows directly from the ultraproduct norm.  Fix
$\varepsilon>0$.  For every $j$, choose $x_j\in E_j$ with $\|x_j\|=1$ and
\[
 \|T_jx_j\|\ge \|T_j\|-\varepsilon.
\]
Then $x=[(x_j)]$ has norm one and
\[
 \|T_\U x\|=\lim_{j\to\U}\|T_jx_j\|
 \ge\lim_{j\to\U}\|T_j\|-\varepsilon.
\]
Letting $\varepsilon\downarrow0$ proves the reverse bound.
\end{proof}

\begin{theorem}\label{thm:main}
Fix $1\le p<\infty$.  Let $I$ be a directed set and
\[
 \big((A_i)_{i\in I},(\varphi_{j,i})_{i\le j}\big)
\]
be a direct system of $L^p$-operator algebras whose connecting homomorphisms
are contractive.  Then its direct limit in the category of Banach algebras and
contractive homomorphisms is an $L^p$-operator algebra.
\end{theorem}

\begin{proof}
For $i\in I$, put $I_i=\{j\in I:j\ge i\}$.  The family of final segments
$(I_i)_{i\in I}$ has the finite-intersection property, so choose an ultrafilter
$\U$ on $I$ which contains every $I_i$.

For each $j\in I$, choose an isometric algebra homomorphism
\[
 \rho_j:A_j\longrightarrow\B(E_j),
 \qquad E_j=L^p(X_j,\mu_j),
\]
with $E_j\ne\{0\}$.  (If necessary, take the $p$-direct sum of a given
representation with the zero representation on $\mathbb C$.)
Let $E=(E_j)_\U$.  The standard $L^p$ ultraproduct theorem says that $E$ is
isometrically an $L^p$-space; see Heinrich~\cite{Heinrich}.

For $i\in I$ and $a\in A_i$, define a uniformly bounded operator family by
\[
 T_j^{i,a}=
 \begin{cases}
  \rho_j(\varphi_{j,i}(a)),&j\ge i,\\
  0,&j\not\ge i.
 \end{cases}
\]
Contractivity gives $\sup_j\|T_j^{i,a}\|\le\|a\|$.  Let
$T^{i,a}\in\B(E)$ be its induced ultraproduct operator.

We use this to define an operator on the dense subalgebra
\[
 A_0=\bigcup_{i\in I}\varphi_i(A_i)\subset A
\]
by setting $\rho(\varphi_i(a))=T^{i,a}$.  We first check that this is well
defined.  Suppose $\varphi_i(a)=\varphi_k(b)$ in $A$.  Choose
$\ell\ge i,k$ and set
\[
 d=\varphi_{\ell,i}(a)-\varphi_{\ell,k}(b)\in A_\ell.
\]
By \eqref{eq:limitnorm},
\[
 \lim_{j\ge\ell}\|\varphi_{j,\ell}(d)\|
 =\|\varphi_\ell(d)\|=0.
\]
On the final segment $I_\ell$, the difference between the two operator
families is $\rho_j(\varphi_{j,\ell}(d))$.  Lemma~\ref{lem:operatornorm}
therefore gives $T^{i,a}=T^{k,b}$.

Linearity and multiplicativity are checked on a common final segment.  For
example, if $\ell\ge i,k$, then for every $j\ge\ell$,
\begin{align*}
 T_j^{i,a}T_j^{k,b}
 &=\rho_j\!\left(\varphi_{j,i}(a)\varphi_{j,k}(b)\right)\\
 &=\rho_j\!\left(
   \varphi_{j,\ell}\big(\varphi_{\ell,i}(a)
                    \varphi_{\ell,k}(b)\big)\right),
\end{align*}
which is the tail family representing the product in $A_0$.

Finally, Lemma~\ref{lem:operatornorm}, the isometry of each $\rho_j$, and
\eqref{eq:limitnorm} give
\begin{align*}
 \|\rho(\varphi_i(a))\|
 &=\lim_{j\to\U}\|T_j^{i,a}\|\\
 &=\lim_{j\to\U}\|\varphi_{j,i}(a)\|\\
 &=\lim_{j\ge i}\|\varphi_{j,i}(a)\|\\
 &=\|\varphi_i(a)\|.
\end{align*}
Thus $\rho:A_0\to\B(E)$ is an isometric algebra homomorphism.  Since $A_0$
is dense in $A$, it extends uniquely to an isometric homomorphism
$A\to\B(E)$.  Its range is closed because $A$ is complete.  Since $E$ is an
$L^p$-space, $A$ is an $L^p$-operator algebra.
\end{proof}

\section{Scope, prior art, and novelty check}

Theorem~\ref{thm:main} answers the page-35 question throughout the finite-$p$
range of the source paper.  It does not assume that the connecting maps are
injective, isometric, nondegenerate, or completely contractive, and the index
set need not be countable.  No claim is made for $p=\infty$.

Phillips's 2014 GPOTS talk summary~\cite[Theorem~4.1]{PhillipsTalk} announces
that a direct limit of $L^p$-operator algebras with \emph{completely
contractive} maps is an $L^p$-operator algebra and mentions the same
ultraproduct mechanism.  The source problem, however, assumes only ordinary
contractivity.  The proof above uses neither matrix norms nor amplifications,
and so removes that additional hypothesis.

The local run indexes were searched for arXiv:1309.6406, the exact title,
``direct limit,'' ``$L^p$-operator algebra,'' and ``ultraproduct.''  Bounded web
searches on August 9, 2026 used the exact source sentence and combinations of
the author, title, contractive/completely contractive maps, direct limits, and
ultraproducts.  They found the source, the 2014 talk's completely-contractive
special case, and later special direct-limit constructions, but no treatment
of arbitrary contractive connecting homomorphisms.  Novelty confidence is
therefore positive but conservative pending expert literature review.

No computation is used.  The proof audit in \path{VERIFICATION.md} identifies
the well-definedness and norm-preservation steps for focused human review.

\begin{thebibliography}{9}
\bibitem{PhillipsSource}
N. C. Phillips,
\emph{Crossed products of $L^p$ operator algebras and the K-theory of Cuntz
algebras on $L^p$ spaces}, arXiv:1309.6406 (2013).

\bibitem{Heinrich}
S. Heinrich,
\emph{Ultraproducts in Banach space theory}, J. Reine Angew. Math.
\textbf{313} (1980), 72--104; doi:10.1515/crll.1980.313.72.

\bibitem{PhillipsTalk}
N. C. Phillips,
\emph{Operator algebras on $L^p$ spaces which ``look like'' $C^*$-algebras},
GPOTS talk summary (2014), Theorem 4.1.
\end{thebibliography}

\end{document}
